% chapters/tool_use.tex — Tool Use chapter (narrative text only)
\chapter{Tool Use and External Interfaces}
\label{chap:tooluse}

Tool-augmented \LLM{}s extend the basic perception--action loop
(Figure~\ref{fig:agent-loop}) with the ability to invoke external
interfaces: calculators, code interpreters, search engines, and
domain-specific APIs~\cite{wu2025agenticreasoningtools}.
\index{tool use}
\index{external interface}

\section{Bounded Tool-Use Policies}
\label{sec:tooluse-policies}

A bounded tool-use policy constrains the agent to a finite, pre-declared
set of tools $\mathcal{T}$, each with a typed schema.
The agent selects a tool $t \in \mathcal{T}$ and supplies structured
arguments; the tool returns an observation that is appended to the agent's
context window.
This keeps the action space tractable and makes the agent's behaviour
auditable~\cite{wu2025agenticreasoningtools}.

\section{Mathematical Reasoning with Tools}
\label{sec:tooluse-math}

\cite{liu2025agenticmath} showed that \LLM{}s equipped with a Python
interpreter as a tool dramatically outperform pure chain-of-thought on
competition-level mathematics.
The agent generates code, executes it, observes the result, and revises
its solution — a verify-and-revise loop that mirrors the iterative
problem-solving strategy of human mathematicians.
\index{mathematical reasoning}

\section{Tool Use and the Planning Objective}
\label{sec:tooluse-planning}

Tool invocations can be modelled as a special class of action in the
planning objective of Equation~\eqref{eq:planning-objective}: a tool call
$a_t \in \mathcal{T}$ transitions the state by adding new information to
the agent's context, yielding a high immediate reward when the retrieved
information resolves an ambiguity or enables a correct next step.
